\section{Virtualização Baseada \textit{Machine Mode}}

Originalmente, o \textit{Machine Mode} (M-mode) foi projetado para ser  o nível de privilégio mais elevado da arquitetura RISC-V, sendo
responsável pela execução do firmware base e pelo controle direto dos recursos de hardware.
Sua função primária consiste na inicialização da plataforma e na configuração dos subsistemas fundamentais antes da
transferência de controle para o \gls{so} hospedeiro.

Dada a natureza restritiva do acesso a determinados recursos de hardware — muitas vezes limitados exclusivamente ao M-mode —, a
arquitetura define camadas de abstração como a \gls{sbi}. Essa interface atua como uma ponte padronizada e segura, permitindo que
níveis de privilégio inferiores solicitem serviços privilegiados ao firmware residente no M-mode.

No contexto deste trabalho, o M-mode é utilizado de forma não convencional. Para isso, deixamos de usar o M-Mode como um simples
\gls{sbi} e passamos a usalo para dar suporte a implementação de um hipervisor completo.

Essa abordagem explora o fato de que o M-mode possui autoridade irrestrita sobre todos os demais níveis de privilégio, permitindo
interceptar, controlar e gerenciar a execução dos sistemas convidados. Para viabilizar esse controle, o hipervisor utiliza
os recursos exclusivos disponibilizados nesse modo, incluindo seu conjunto dedicado de \glspl{csr}, além de mecanismos próprios
para tratamento de interrupções e gerenciamento de temporização. Dessa forma, o M-mode fornece a soberania necessária
para reproduzir, por software, e as vezes até em hardware, funcionalidades que normalmente só seriam oferecidas por
hardwares especializados compatíveis com as extensões de virtualização.
